\chapter{Analyse en Requirements}
\label{ch:analyse}

In dit hoofdstuk worden de functionele en technische vereisten van de Proof of Concept (PoC) gedefinieerd. 
Waar de literatuurstudie in Hoofdstuk 2 de theoretische noodzaak voor Mixed Reality (MR) binnen de zorg vaststelde, 
vertaalt dit hoofdstuk die noodzaak naar concrete systeemeisen. De focus ligt op het overbruggen van de kloof tussen de fysieke hardware van de 
reanimatiepop en de digitale omgeving van de Meta Quest 3.

\section{Probleemcontext}
De huidige reanimatietrainingen lijden onder het zogenaamde \textit{split-attention effect}, waarbij de cursist de aandacht moet verdelen tussen de 
pop en een extern feedbackscherm. De voorgestelde oplossing beoogt deze cognitieve belasting te verlagen door data direct in het gezichtsveld van de 
verpleegkundige te projecteren. De technische uitdaging is echter dat de gebruikte Laerdal Little Anne QCPR pop een gesloten systeem is zonder publiek toegankelijke API.

\section{Functionele Requirements}
De functionele eisen beschrijven wat de applicatie moet kunnen om als een effectieve trainingstool te dienen.

\begin{itemize}
    \item \textbf{Draadloze Data-acquisitie:} De applicatie moet via Bluetooth Low Energy (BLE) verbinding maken met de reanimatiepop.
    \item \textbf{Real-time Feedback-lus:} Elke compressie moet onmiddellijk resulteren in een visuele verandering in de MR-omgeving.
    \item \textbf{Visuele Indicatoren:} Het systeem moet conform de ERC 2025 richtlijnen aangeven of de compressiediepte (5--6 cm) en de frequentie (100--120 bpm) correct zijn.
    \item \textbf{Ghost Avatar:} Een transparante referentie-avatar moet het ideale bewegingspatroon demonstreren om de student te begeleiden (\textit{feedforward}).
\end{itemize}

\section{Niet-functionele Requirements}
De niet-functionele eisen bepalen de kwaliteit en de bruikbaarheid van de software.

\begin{itemize}
    \item \textbf{Robuustheid:} De BLE-verbinding moet stabiel blijven in een omgeving met potentieel veel interferentie (zoals een ziekenhuis- of klasomgeving).
    \item \textbf{Intuïtieve UI:} De interface moet de \textit{situational awareness} van de verpleegkundige versterken in plaats van deze te belemmeren met complexe menu's.
\end{itemize}

\section{Technische Randvoorwaarden}
Voor de ontwikkeling van deze PoC zijn de volgende hardware- en softwarekeuzes gemaakt:

\begin{description}
    \item[Meta Quest 3:] Gekozen vanwege de hoge-resolutie \textit{video passthrough} en de krachtige \textit{spatial computing} mogelijkheden voor hand-tracking.
    \item[Unity Engine:] Gebruikt als primaire ontwikkelomgeving vanwege de volwassenheid van de Meta XR SDK en de flexibiliteit in het afhandelen van ruwe datastromen.
    \item[Android-gebaseerd OS:] Aangezien de Quest 3 op Android draait, moet de communicatielaag compatibel zijn met de Android Bluetooth-stack.
\end{description}

\section{MoSCoW-Prioritering}
Gezien de beperkte tijdspanne van de bachelorproef en de complexiteit van de hardware-integratie, is de scope vastgelegd middels de MoSCoW-methode:

\begin{table}[ht]
    \centering
    \begin{tabular}{|l|p{10cm}|}
        \hline
        \textbf{Prioriteit} & \textbf{Eis} \\ \hline
        \textbf{Must-have} & Stabiele BLE-datastroom, visualisatie van diepte en tempo, correcte ruimtelijke verankering op de pop. \\ \hline
        \textbf{Should-have} & Implementatie van de Ghost Avatar, kleurcodering bij foutieve handelingen. \\ \hline
        \textbf{Could-have} & Loggen van sessiedata voor latere evaluatie, ondersteuning voor verschillende reanimatie-scenario's. \\ \hline
        \textbf{Won't-have} & Volledige integratie met ziekenhuis-informatiesystemen of multi-player functionaliteit. \\ \hline
    \end{tabular}
    \caption{MoSCoW-prioritering voor de Proof of Concept.}
    \label{tab:moscow_analyse}
\end{table}

\section{Conclusie van de analyse}
De analyse toont aan dat de grootste bottleneck niet de visualisatie in MR is, maar de acquisitie van de live data. Zonder een gedocumenteerde API van de fabrikant is 
een standaard implementatie onmogelijk. Dit noodzaakt een diepgaand onderzoek naar het communicatieprotocol van de pop, wat in het volgende hoofdstuk (Reverse Engineering) 
technisch zal worden uitgewerkt.